\documentclass{article}
\usepackage{multicol}
\title{Multicolumn with Column Breaks}
\author{Ferritex Test}
\date{April 2026}

\begin{document}
\maketitle

\section{Two Columns with Manual Break}

\begin{multicols}{2}
This text appears in the left column. It contains enough words to
demonstrate that the column break directive forces a transition to the
next column regardless of how much vertical space remains.

\columnbreak

This text appears in the right column after the explicit column break.
The break ensures predictable placement without relying on automatic
balancing.
\end{multicols}

\section{Three Columns}

\begin{multicols}{3}
First column content that fills the narrow measure of a three-column
layout.

\columnbreak

Second column with a moderate amount of text to show the middle
position.

\columnbreak

Third and final column wrapping up the three-column demonstration.
\end{multicols}

\end{document}
